\chapter{Technical Architecture}

\section{Overview}

The technical architecture of \textit{ECHO: The Living World} is designed to provide a scalable, modular, and maintainable foundation for developing a modern open-world action role-playing game. The architecture separates game systems into independent modules that communicate through well-defined interfaces, making future expansion and maintenance easier.

The design emphasizes performance, flexibility, and code reusability while supporting the game's large open world, intelligent AI systems, dynamic environments, and story-driven gameplay.

\section{Architecture Layers}

The game architecture is divided into the following major layers:

\begin{itemize}
    \item Presentation Layer
    \item Gameplay Layer
    \item Artificial Intelligence Layer
    \item Physics Layer
    \item Data Management Layer
    \item Audio Layer
    \item Rendering Layer
\end{itemize}

Each layer performs specialized responsibilities while interacting with the others through controlled communication.

\section{Core Gameplay Systems}

The gameplay layer manages the primary mechanics of the game, including:

\begin{itemize}
    \item Player movement
    \item Combat mechanics
    \item Companion interactions
    \item Quest management
    \item Dialogue system
    \item Inventory management
    \item Skill progression
    \item Reputation system
\end{itemize}

These systems work together to provide a seamless gameplay experience.

\section{Artificial Intelligence System}

The AI subsystem controls:

\begin{itemize}
    \item Civilian behaviour
    \item Enemy behaviour
    \item Companion behaviour
    \item Wildlife simulation
    \item Guard patrols
    \item Boss encounters
\end{itemize}

Each AI entity follows behaviour rules based on its role, current environment, and interaction with the player.

\section{Data Management}

Game data is organized into modular assets including:

\begin{itemize}
    \item Character data
    \item Quest data
    \item World data
    \item Dialogue data
    \item Item database
    \item Save files
\end{itemize}

This modular approach simplifies updates and future content additions.

\section{Rendering System}

The rendering system is responsible for:

\begin{itemize}
    \item Lighting
    \item Shadows
    \item Weather effects
    \item Environmental rendering
    \item Character animation
    \item Visual effects
\end{itemize}

The objective is to create a visually immersive world while maintaining stable performance.

\section{Audio System}

The audio subsystem manages:

\begin{itemize}
    \item Background music
    \item Environmental ambience
    \item Character dialogue
    \item Combat sounds
    \item Wildlife sounds
    \item Dynamic sound effects
\end{itemize}

Audio adapts dynamically according to gameplay situations and player actions.

\section{Save System}

The save system records:

\begin{itemize}
    \item Story progression
    \item Player choices
    \item Inventory
    \item Character progression
    \item Companion relationships
    \item Quest completion
    \item World state
\end{itemize}

This ensures that player decisions have persistent consequences throughout the game.

\section{Future Scalability}

The modular architecture supports future improvements such as:

\begin{itemize}
    \item Downloadable content
    \item Additional continents
    \item New quests
    \item New AI behaviours
    \item Multiplayer support
\end{itemize}

\section{Summary}

The technical architecture of \textit{ECHO: The Living World} provides a robust and scalable framework capable of supporting an immersive open-world RPG while allowing future expansion without major redesign.